\documentclass[12pt,a4paper]{article}

\usepackage[a4paper,margin=1in]{geometry}
\usepackage{graphicx}
\usepackage{booktabs}
\usepackage{amsmath,amssymb}
\usepackage{float}
\usepackage{hyperref}
\usepackage{xcolor}
\usepackage{listings}
\usepackage{enumitem}
\usepackage{fancyhdr}
\usepackage{longtable}

\hypersetup{colorlinks=true,linkcolor=blue,urlcolor=blue}

\pagestyle{fancy}
\fancyhf{}
\lhead{ICS1512}
\rhead{Machine Learning Algorithms Laboratory}
\cfoot{\thepage}

\lstset{
language=Python,
basicstyle=\ttfamily\small,
numbers=left,
frame=single,
breaklines=true,
showstringspaces=false
}

\begin{document}

\begin{tabular}{|p{4cm}|p{11cm}|}
\hline
Experiment No. & 6 \\ \hline
Experiment Title & Dimensionality Reduction and Model Evaluation (With and Without PCA)
\footnote{\textbf{GitHub Repository: }\url{https://github.com/username/repository}}\\ \hline
Student Name & \\ \hline
Register Number & \\ \hline
Faculty & Dr. Poreddy Ajay Kumar Reddy \\ \hline
Submission Date & \\ \hline
\end{tabular}

\section{Objective}
To study how Principal Component Analysis (PCA) affects the performance of different classifiers.
Ten models are trained twice each -- once on the original standardized features and once on the
PCA-reduced features -- and the two settings are compared after tuning each model with 5-fold
cross-validation.

\section{Problem Statement}
Given a set of correlated numerical features describing a tumour, the task is to classify it as
malignant or benign. Since several of the original features carry overlapping information, PCA is
applied to compress them into a smaller set of uncorrelated components, and the question this
experiment answers is whether that compression helps, hurts, or makes no real difference to
classification performance, and whether the effect depends on the type of model used.

\section{Dataset Description}
\begin{longtable}{|p{6cm}|p{8cm}|}
\hline
Dataset Name & Wisconsin Diagnostic Breast Cancer Dataset \\ \hline
Dataset Source & UCI Machine Learning Repository (loaded via \texttt{sklearn.datasets.load\_breast\_cancer}) \\ \hline
Number of Samples & 569 \\ \hline
Number of Features & 30 numerical features (mean, standard error, and worst value of 10 cell nucleus measurements) \\ \hline
Number of Classes & 2 (Malignant, Benign) \\ \hline
Missing Values & None \\ \hline
Train-Test Split & 80\% train (455 samples) / 20\% test (114 samples), stratified \\ \hline
\end{longtable}

\section{Exploratory Data Analysis}
Before reducing the feature space it is worth checking why PCA might help: the 30 features are
mostly different summary statistics (mean, error, worst) of the same 10 underlying measurements,
so many of them are strongly correlated with each other.
\begin{itemize}
\item Dataset overview -- 569 rows, no missing values, all features numeric.
\item Class distribution -- roughly 63\% benign and 37\% malignant, a mild but not severe imbalance.
\item Correlation matrix -- a large block of features (radius, perimeter, area) are almost redundant with one another.
\item Feature distribution -- most features are right-skewed with a long tail of larger, more irregular tumours.
\end{itemize}

\subsection*{Figure 1: 12-Subplot Consolidated Exploratory Data Analysis}

\begin{figure}[H]
\centering
\includegraphics[width=\linewidth]{exp6_eda_12_subplots.eps}
\caption{Consolidated 12-subplot exploratory data analysis covering target class balance, missing values, duplicate distribution, correlation heatmap, and key feature distributions}
\end{figure}

\textbf{Inference:}

The dataset is moderately imbalanced towards the benign class, so accuracy alone is not enough
and precision, recall, and F1-score are tracked as well. The correlation heatmap shows large
red blocks, confirming that several groups of features move together almost perfectly. This is
exactly the kind of redundancy PCA is designed to compress, since keeping 30 near-duplicate axes
adds little information but does add noise and computation.

\section{Data Preprocessing}

\begin{itemize}
\item \textbf{Missing value handling:} not required, the dataset has no missing entries.
\item \textbf{Encoding:} the target is already binary (0 = malignant, 1 = benign), no encoding needed.
\item \textbf{Feature scaling:} all 30 features are standardized (zero mean, unit variance) using
\texttt{StandardScaler}, fit only on the training set and then applied to the test set, since
several of the downstream models (SVM, KNN, Logistic Regression) are sensitive to feature scale.
\item \textbf{Train-test split:} 80--20 stratified split so both sets keep the same class ratio.
\item \textbf{Dimensionality reduction:} PCA is fit on the standardized training data only. Instead
of picking an arbitrary number of components, the number of components needed to explain 95\% of
the variance was used, which came out to 10 components (out of 30), retaining 95.27\% of the
variance.
\end{itemize}

\subsection*{Figure 2: PCA Scree Plot}

\begin{figure}[H]
\centering
\includegraphics[width=0.75\linewidth]{fig_pca_scree.eps}
\caption{Cumulative explained variance vs number of principal components}
\end{figure}

\textbf{Inference:}

The curve rises steeply for the first few components and then flattens out, which is consistent
with the correlation heatmap -- a handful of components already capture most of the spread in the
data. Ten components is a reasonable cut-off: it keeps two-thirds of the original dimensionality
reduction (10 out of 30 features) while still retaining 95\% of the information.

\begin{longtable}{|p{6cm}|p{8cm}|}
\hline
\textbf{PCA Summary} & \\ \hline
Variance Target & 95\% \\ \hline
Components Chosen & 10 (out of 30) \\ \hline
Explained Variance Achieved & 95.27\% \\ \hline
\end{longtable}

\section{Machine Learning Algorithm}
Ten classifiers were compared, spanning linear, margin-based, instance-based, and ensemble methods:

\begin{itemize}
\item \textbf{Logistic Regression / SVM:} both find a decision boundary in feature space (linear
or maximum-margin respectively); their working principle relies directly on distances/dot products
in feature space, so they are the models most likely to be affected by changing that space via PCA.
\item \textbf{Naive Bayes:} assumes features are conditionally independent given the class and
estimates a Gaussian distribution per feature; PCA's uncorrelated components actually fit this
assumption better than the original correlated features do, in principle.
\item \textbf{KNN:} classifies a point by the majority label among its nearest neighbours in
feature space, so it is directly affected by which axes are kept and how they are scaled.
\item \textbf{Decision Tree / Random Forest / AdaBoost / Gradient Boosting / XGBoost:} all split
the feature space one axis (or one linear combination, for boosting variants) at a time based on
impurity reduction; trees can already ignore uninformative features on their own, so they are
expected to gain less from PCA than the distance-based models.
\item \textbf{Stacking:} combines SVM, Naive Bayes, and a Decision Tree as base learners, with a
Logistic Regression meta-learner that learns how much to trust each base model's prediction.
\end{itemize}

\textbf{Advantages of PCA:} removes redundant/correlated information, speeds up training, can
reduce overfitting for high-variance models.

\textbf{Limitations of PCA:} components are linear combinations of the original features and lose
direct interpretability; if the classes are not separated along directions of high variance, PCA
can discard exactly the information a classifier needed.

\section{Experimental Setup}
\begin{tabular}{ll}
\toprule
Python Version & 3.x \\
Scikit-Learn Version & 1.8 \\
IDE & Jupyter Notebook \\
Operating System & \\
Hardware & \\
\bottomrule
\end{tabular}

\section{Hyperparameter Selection}
Each model was tuned with \texttt{GridSearchCV} (5-fold, accuracy scoring) separately for the
No-PCA and With-PCA settings. The best parameters found for each are below.

\begin{longtable}{p{3.2cm}p{2cm}p{5.5cm}p{2cm}}
\toprule
Model & Setting & Best Parameters & Best CV Accuracy \\
\midrule
SVM & No-PCA & C=0.1, kernel=linear, gamma=scale & 0.9758 \\
SVM & With-PCA & C=1, kernel=linear, gamma=scale & 0.9780 \\
Naive Bayes & No-PCA & var\_smoothing=1e-9 & 0.9341 \\
Naive Bayes & With-PCA & var\_smoothing=1e-9 & 0.9165 \\
KNN & No-PCA & n\_neighbors=9, weights=uniform & 0.9670 \\
KNN & With-PCA & n\_neighbors=5, weights=uniform & 0.9648 \\
Logistic Regression & No-PCA & C=0.1, penalty=l2 & 0.9824 \\
Logistic Regression & With-PCA & C=0.1, penalty=l2 & 0.9824 \\
Decision Tree & No-PCA & criterion=entropy, max\_depth=3 & 0.9319 \\
Decision Tree & With-PCA & criterion=entropy, max\_depth=3 & 0.9297 \\
Random Forest & No-PCA & n\_estimators=100, max\_depth=10 & 0.9626 \\
Random Forest & With-PCA & n\_estimators=200, max\_depth=None & 0.9604 \\
AdaBoost & No-PCA & n\_estimators=100, learning\_rate=1.0 & 0.9802 \\
AdaBoost & With-PCA & n\_estimators=100, learning\_rate=0.5 & 0.9648 \\
Gradient Boosting & No-PCA & n\_estimators=150, max\_depth=2, lr=0.1 & 0.9648 \\
Gradient Boosting & With-PCA & n\_estimators=150, max\_depth=2, lr=0.1 & 0.9604 \\
XGBoost & No-PCA & n\_estimators=100, max\_depth=3, lr=0.1 & 0.9714 \\
XGBoost & With-PCA & n\_estimators=100, max\_depth=5, lr=0.05 & 0.9648 \\
Stacking & No-PCA & final\_estimator C=10 & 0.9692 \\
Stacking & With-PCA & final\_estimator C=10 & 0.9736 \\
\bottomrule
\end{longtable}

\section{Results}
Test-set performance for all ten models, No-PCA vs With-PCA.

\begin{longtable}{p{3.2cm}p{1.6cm}p{1.4cm}p{1.6cm}p{1.3cm}p{1.4cm}p{1.4cm}}
\toprule
Model & Setting & Accuracy & Precision & Recall & F1 & ROC-AUC \\
\midrule
SVM & No-PCA & 0.9825 & 0.9861 & 0.9861 & 0.9861 & 0.9937 \\
SVM & With-PCA & 0.9649 & 0.9857 & 0.9583 & 0.9718 & 0.9950 \\
Naive Bayes & No-PCA & 0.9298 & 0.9444 & 0.9444 & 0.9444 & 0.9868 \\
Naive Bayes & With-PCA & 0.9211 & 0.9315 & 0.9444 & 0.9379 & 0.9709 \\
KNN & No-PCA & 0.9737 & 0.9600 & 1.0000 & 0.9796 & 0.9944 \\
KNN & With-PCA & 0.9561 & 0.9589 & 0.9722 & 0.9655 & 0.9788 \\
Logistic Regression & No-PCA & 0.9737 & 0.9726 & 0.9861 & 0.9793 & 0.9957 \\
Logistic Regression & With-PCA & 0.9737 & 0.9726 & 0.9861 & 0.9793 & 0.9954 \\
Decision Tree & No-PCA & 0.9474 & 0.9459 & 0.9722 & 0.9589 & 0.9448 \\
Decision Tree & With-PCA & 0.9386 & 0.9577 & 0.9444 & 0.9510 & 0.9559 \\
Random Forest & No-PCA & 0.9561 & 0.9589 & 0.9722 & 0.9655 & 0.9939 \\
Random Forest & With-PCA & 0.9386 & 0.9577 & 0.9444 & 0.9510 & 0.9864 \\
AdaBoost & No-PCA & 0.9561 & 0.9467 & 0.9861 & 0.9660 & 0.9818 \\
AdaBoost & With-PCA & 0.9386 & 0.9577 & 0.9444 & 0.9510 & 0.9854 \\
Gradient Boosting & No-PCA & 0.9474 & 0.9459 & 0.9722 & 0.9589 & 0.9911 \\
Gradient Boosting & With-PCA & 0.9298 & 0.9444 & 0.9444 & 0.9444 & 0.9854 \\
XGBoost & No-PCA & 0.9474 & 0.9459 & 0.9722 & 0.9589 & 0.9934 \\
XGBoost & With-PCA & 0.9386 & 0.9577 & 0.9444 & 0.9510 & 0.9931 \\
Stacking & No-PCA & 0.9825 & 0.9861 & 0.9861 & 0.9861 & 0.9954 \\
Stacking & With-PCA & 0.9474 & 0.9714 & 0.9444 & 0.9577 & 0.9924 \\
\bottomrule
\end{longtable}

\section{Result Visualization}

\subsection*{Figure 3: Test Accuracy, No-PCA vs With-PCA}
\begin{figure}[H]
\centering
\includegraphics[width=0.85\linewidth]{fig_accuracy_comparison.eps}
\caption{Test accuracy of all ten models under both settings}
\end{figure}
\textbf{Inference:} SVM and Stacking are the strongest models overall, both reaching 98.25\%
accuracy without PCA. Almost every model's bar drops slightly under With-PCA, the only two that
do not lose ground are Logistic Regression (unchanged) and SVM's cross-validation accuracy
(marginally higher, though its test accuracy still drops). This suggests PCA is, at best, neutral
and at worst a small cost on this dataset.

\subsection*{Figure 4: 5-Fold CV Stability, No-PCA vs With-PCA}
\begin{figure}[H]
\centering
\includegraphics[width=0.85\linewidth]{fig_cv_stability.eps}
\caption{Mean 5-fold cross-validation accuracy per model, both settings}
\end{figure}
\textbf{Inference:} Stacking and SVM are the only two models whose CV accuracy actually improves
with PCA (+0.44\% and +0.22\% respectively). Naive Bayes loses the most (-1.76\%), followed by
AdaBoost (-1.54\%) and XGBoost (-0.66\%). Tree-based ensembles consistently lose a small amount
of accuracy under PCA, supporting the idea that they were already handling the original
correlated features fine on their own.

\subsection*{Figure 5: ROC Curves}
\begin{figure}[H]
\centering
\includegraphics[width=0.75\linewidth]{fig_roc_curves.eps}
\caption{ROC curves for all models, No-PCA setting}
\end{figure}
\textbf{Inference:} nearly all curves hug the top-left corner with AUC above 0.98, meaning every
model separates the two classes well regardless of PCA; ROC-AUC is far less sensitive to PCA than
raw accuracy is, since the ranking of predicted probabilities barely changes even when a few
borderline points flip class near the 0.5 threshold.

\subsection*{Figure 6: Confusion Matrix -- Best Model}
\begin{figure}[H]
\centering
\includegraphics[width=0.6\linewidth]{fig_confusion_matrix.eps}
\caption{Confusion matrix of the best-performing model (SVM / Stacking, No-PCA)}
\end{figure}
\textbf{Inference:} the few errors that remain are almost all malignant cases predicted as
benign (false negatives), which matters more in a medical context than the reverse, and is worth
flagging even though overall accuracy is high.

\section{Discussion}
PCA compressed 30 features down to 10 while keeping 95\% of the variance, which is a meaningful
reduction in dimensionality. However, it did not clearly improve classification performance on
this dataset: only Stacking and SVM saw a (small) gain in cross-validation accuracy, while the
tree-based ensembles (Random Forest, AdaBoost, Gradient Boosting, XGBoost) and Naive Bayes all
performed slightly worse with PCA applied. This matches the theory -- linear and margin-based
models benefit most from having redundant, correlated features collapsed into a smaller
orthogonal set, while tree-based models already select the useful features themselves at each
split and do not need PCA to do that for them.

Logistic Regression is a useful sanity check here: its best hyperparameters and CV accuracy are
identical in both settings, meaning the extra components dropped by PCA were not adding meaningful
signal, and it was just as easy to overfit or underfit either way.

A clear limitation of this experiment is the small grid used for hyperparameter tuning (2--4
options per parameter), which was a deliberate trade-off to keep 10 models $\times$ 2 settings
$\times$ 5-fold CV runnable in reasonable time; a wider grid could shift the exact numbers, though
the broad pattern (PCA helping linear/margin models, not helping trees) is unlikely to change.

\section{Conclusion}
PCA is not a universal improvement -- its effect depends heavily on the type of model. On this
dataset, SVM and Stacking mildly benefited from PCA in cross-validation, Logistic Regression was
unaffected, and every tree-based ensemble along with Naive Bayes performed slightly worse. The
best overall test accuracy (98.25\%) was achieved without PCA, by SVM and Stacking. PCA remains
useful when training speed or feature redundancy is a real concern, but for this dataset and this
set of models it is closer to a neutral trade-off than a clear win.

\section{References}
\begin{enumerate}
\item Scikit-learn Developers, ``sklearn.decomposition.PCA,'' scikit-learn documentation. [Online]. Available: https://scikit-learn.org/stable/modules/generated/sklearn.decomposition.PCA.html
\item Scikit-learn Developers, ``Supervised learning,'' scikit-learn documentation. [Online]. Available: https://scikit-learn.org/stable/supervised\_learning.html
\item Scikit-learn Developers, ``Model selection and evaluation,'' scikit-learn documentation. [Online]. Available: https://scikit-learn.org/stable/model\_selection.html
\item W. Wolberg, O. Mangasarian, N. Street, and W. Street, ``Breast Cancer Wisconsin (Diagnostic) Data Set,'' UCI Machine Learning Repository. [Online]. Available: https://archive.ics.uci.edu
\end{enumerate}

\end{document}
